\chapter*{Abstract}
\addcontentsline{toc}{chapter}{Abstract}

This thesis asks whether longitudinal transitions in fuzzy cluster membership have incremental explanatory power for changes in the market value of football players, beyond individual performance statistics. It applies the robust fuzzy $C$-medoids model for mixed data with a noise cluster to eight seasons (2017/18 to 2024/25) of the Portuguese Primeira Liga, aligns the archetypes across seasons, and relates the change in each player's membership vector to the change in his Transfermarkt value for 1,361 season pairs. A two-stage bootstrap accounts for the estimation of the memberships.

The clustering yields four interpretable archetypes (defenders, central midfielders, wide attackers and strikers) and a small noise cluster; 75\% of the players keep their archetype from one season to the next. In the main configuration, membership changes are associated with changes in value beyond the baselines (bootstrap $p=0.015$ against the full baseline and $0.079$ against the core baseline), with gains of 0.3 to 0.6 percentage points of adjusted $R^2$. The association is weak. It depends on the number of archetypes and on the random draw of the clustering, no result survives a correction for multiple comparisons, and it largely disappears when changes in the position labels, which enter the clustering but not the baselines, are controlled (bootstrap $p=0.35$ and $0.12$).

The thesis therefore finds no robust incremental explanatory power of membership transitions for changes in market value. Its main contributions are the description of the archetypes and of their dynamics, and a documented analysis of how sensitive the results of a fuzzy clustering are to its tuning, to the random draw, and to the controls.

\bigskip\noindent\textbf{Keywords:} fuzzy clustering, noise cluster, football analytics, market value, Primeira Liga
\clearpage
